\exercisegroup

\begin{enumerate}
\def\labelenumi{\arabic{enumi})}
\tightlist
\item
  设 \(\mathbf{E}\) 是非离散拓扑域 \(\mathbf{K}\) 上的一个左拓扑向量空间。称 \(\mathbf{B}\) （\(\mathbf{E}\) 的子集）是有界的，如果对 0 的每个邻域 \(\mathbf{V}\) （\(\mathbf{E}\) 中），存在 \(\lambda \neq 0\) （\(\mathbf{K}\) 中），使得 \(\lambda \mathbf{B} \subset \mathbf{V}\) 。
\end{enumerate}

\begin{enumerate}
\def\labelenumi{\alph{enumi})}
\item
  证明：若 B 在 E 中有界，则对 E 中 0 的每个邻域 V，存在 K 中 0 的一个邻域 U，使得 U.B ⊂ V。
\item
  证明：E 中有界集的闭包是有界的。两个有界集的并是有界的。E 中每个预紧集都是有界的。把 III, 第 4 页, 命题 4 的诸推论推广到 K 上的拓扑向量空间。
\item
  证明：若 A 是 K 中的有界集（把 K 视为自身上的左向量空间），B 是 E 中的有界集，则 A.B 在 E 中有界。
\item
  把 III, 第 4 页 的命题 3 推广到 K 是可度量化拓扑除环的情形。e) 把拟完备空间的概念及其性质推广到拓扑向量空间。
\end{enumerate}

\begin{enumerate}
\def\labelenumi{\arabic{enumi})}
\setcounter{enumi}{1}
\tightlist
\item
  \begin{enumerate}
  \def\labelenumii{\alph{enumii})}
  \tightlist
  \item
    设 E 是非离散拓扑域 K 上的一个左拓扑向量空间。证明：若存在 E 中 0 的一个有界邻域 V（习题 1），则诸集 \(\lambda\) V（\(\lambda \in \mathbf{K}\) 且 \(\lambda \neq 0\) ）构成 E 中 0 的一个基本邻域系。若 K 可度量化，则与 E 的拓扑相伴的 Hausdorff 拓扑（GT, III, § 2, No.~6）可度量化。若 \(\mathbf{K} = \mathbf{R}\) 或 \(\mathbf{K} = \mathbf{C}\) ，则 E 上比 E 的给定拓扑更粗的诸拓扑之中的最细局部凸拓扑（II, 第 80 页, 习题 23）可由单个半范数定义。
  \end{enumerate}
\end{enumerate}

\begin{enumerate}
\def\labelenumi{\alph{enumi})}
\setcounter{enumi}{1}
\item
  证明：局部凸 Hausdorff 空间（其中没有一个仅是点 0）的无穷乘积的拓扑不能由单个半范数定义。
\item
  设 E 是一个局部凸空间，其拓扑由半范数的递增序列 \((p_n)\) 定义。要使 E 的拓扑可由单个半范数定义，必须且只需存在整数 \(n_0\) ，使得对每个 \(n \geqslant n_0\) ，存在数 \(k_n \geqslant 0\) ，使得 \(p_n(x) \leqslant k_n p_{n_0}(x)\) 对一切 \(x \in E\) 成立。
\end{enumerate}

\(d\) ) 设 E 是 \(\mathbf{R}\) 上的向量空间，由区间 \(\mathrm{I} = [0,1]\) 上无穷可微的数值函数组成。对每个整数 \(n\geqslant 0\) ，令 \(p_n(f) = \sup_{0\leqslant k\leqslant n}\left(\sup_{x\in I}|f^{(k)}(x)|\right)\) （其中

\(f^{(0)} = f)\) ）；证明诸 \(p_n\) 是 E 上的范数，且由范数序列 \(p_n\) 定义的拓扑不能由单个范数定义。

\begin{enumerate}
\def\labelenumi{\arabic{enumi})}
\setcounter{enumi}{2}
\item
  设 E 是 R 上的可度量化向量空间，d 是与 E 的拓扑相容的平移不变距离。令 \(|x| = d(x, 0)\) （I, 第 16 页）。证明：对每个整数 n \textgreater{} 0，有 \(\frac{1}{n}|x| \leqslant \left|\frac{x}{n}\right|\) 。由此推出：若 B 是 E 的有界子集，则 \(\sup_{x \in B} |x| < +\infty\) （换言之，B 对距离 d 有界（GT, IX, § 2, No.~3））。给出一个可度量化向量空间 E 以及 E 中一个集合的例子，该集合无界但对距离 d 有界（习题 2）。
\item
  设 E 是非离散可度量化拓扑域 K 上的一个拓扑向量空间。证明：若 E 是 Baire 空间，且 E 中由有界子集构成的有界结构具有可数基（III, 第 37 页, 习题 1），则存在 E 中 0 的一个有界邻域，从而与 E 的拓扑相伴的 Hausdorff 拓扑可度量化（III, 第 37 页, 习题 2）（与习题 6 比较）。
\item
  设 \(\mathbf{E}\) 是非离散赋值域 \(\mathbf{K}\) 上的可度量化向量空间。证明：若 \((\mathbf{B}_n)\) 是 \(\mathbf{E}\) 的有界子集的任意一个序列（III, 第 37 页, 习题 1），则存在标量的一个序列 \((\lambda_n)\) （诸标量 \(\neq 0\)），使得诸集 \(\lambda_n\mathbf{B}_n\) 的并是有界的。
\item
  设 \(\mathrm{E}\) 是一个局部凸空间，它是局部凸 Hausdorff 空间的严格递增序列 \((\mathrm{E}_n)\) 的严格归纳极限，且每个 \(\mathrm{E}_n\) 在 \(\mathrm{E}_{n+1}\) 中闭（II, 第 32 页, 命题 9）。
\end{enumerate}

\begin{enumerate}
\def\labelenumi{\alph{enumi})}
\tightlist
\item
  证明 E 不可度量化（利用 III, 第 5 页, 命题 6 以及前面的习题 5）。\\
\item
  要使 E 的典范有界结构具有可数基，必须且只需每个 \(\mathrm{E}_n\) 的典范有界结构在 \(\mathrm{E}_n\) 中具有可数基。
\end{enumerate}

\begin{enumerate}
\def\labelenumi{\arabic{enumi})}
\setcounter{enumi}{6}
\tightlist
\item
  \begin{enumerate}
  \def\labelenumii{\alph{enumii})}
  \tightlist
  \item
    设 E 是一个无穷维 Banach 空间，S 是 E 的紧凸均衡子集所成的族，它对关系 \(\subset\) 是有向集。证明：E 是 Banach 空间 \(E_{A}\) （A 取遍 S）的归纳系统的归纳极限。（用反证法证明：对诸 \(E_{A}\) 的拓扑的归纳极限拓扑而言，0 的邻域 V 含有一个以 0 为中心的球；为此，注意否则将存在 E 中的点列 \((x_{n})\) ，使得 \(\|x_{n}\| \leqslant 1/n^{2}\) ，且不属于 V。）由此推出：E 中存在不含于任何 \(E_{A}\) （\(A \subset S\) ）的有界子集。
  \end{enumerate}
\end{enumerate}

\begin{enumerate}
\def\labelenumi{\alph{enumi})}
\setcounter{enumi}{1}
\tightlist
\item
  设 E 是一个无穷维 Banach 空间。在向量空间 \(\prod_{m=1}^{\infty}E_{m}\) （其中对每个 m 有 \(E_{m}=E\) ）上，用 \(T_{n}\) 表示这样的拓扑：在每个 \(E_{m}\) （\(m\leqslant n\) ）上取 Banach 空间拓扑的乘积，而对 m\textgreater n 取 \(E_{m}\) 上的最细局部凸拓扑；用 \(F_{n}\) 表示局部凸空间 \(\prod_{m=1}^{\infty}E_{m}\) （赋以拓扑 \(T_{n}\) ）。每个恒等映射 \(F_{n}\to F_{n+1}\) 都连续；证明：归纳系统 \(F_{n}\) 的归纳极限空间是空间 F，其中 F 是在 \(\prod_{m=1}^{\infty}E_{m}\) 上赋以各因子的 Banach 空间拓扑之乘积而得到的空间。由此推出：F 中存在在任何 \(F_{n}\) 中都无界的有界子集。
\end{enumerate}

\begin{enumerate}
\def\labelenumi{\arabic{enumi})}
\setcounter{enumi}{7}
\item
  证明：在一个由拓扑向量空间（在 \(\mathbf{R}\) 或 \(\mathbf{C}\) 上，没有一个仅是点 0）作成的无穷乘积的空间中，典范有界结构不具有可数基（首先证明只需对空间 \(\mathbf{R}^{\mathbf{N}}\) 证明这一点，然后利用 III, 第 38 页, 习题 4）。
\item
  设 E 是 R 上由区间 I = \{0, 1\} 上全体调节函数（FVR, I, 第 4 页）组成的向量空间。对每个整数 n \textgreater{} 0，令 \(V_{n}\) 为全体函数 \(f \in E\) （满足 \(\int_{0}^{1} \sqrt{|f(t)|} dt \leqslant 1/n\) ）所成的集。证明：诸集 \(V_{n}\) 构成 E 上一个与 E 的向量空间结构相容的可度量化拓扑的 0 的基本邻域系，且对该拓扑诸集 \(V_{n}\) 是有界的；但每个 \(V_{n}\) 的凸包是整个空间 E。（注意每个函数 \(f \in E\) 都可写成 \(f = \frac{1}{2}(g + h)\) ，其中 g 与 h 属于 E，且
\end{enumerate}

\[
\int_ {0} ^ {1} \sqrt {| g (t) |} d t = \int_ {0} ^ {1} \sqrt {| h (t) |} d t = \frac {1}{\sqrt {2}} \int_ {0} ^ {1} \sqrt {| f (t) |} d t.
\]

由此推出：E 上比 E 的拓扑更粗的唯一的局部凸拓扑是 E 上的最粗拓扑。

\begin{enumerate}
\def\labelenumi{\arabic{enumi})}
\setcounter{enumi}{9}
\tightlist
\item
  设 \((\mathrm{E}_{\iota})_{\iota \in I}\) 是非离散拓扑域 K 上由 Hausdorff 拓扑向量空间组成的一个无穷族，其中没有一个仅是点 0。设 E 是诸 \(\mathbf{E}_{\iota}\) 的直和向量空间，\(\mathcal{T}_0\) 是 I, 第 24 页, 习题 14 中定义的 E 上的拓扑。那么，E 的子集 B 对 \(\mathcal{T}_0\) 有界（III, 第 37 页, 习题 1），当且仅当 B 含于某个乘积子空间 \(\prod_{\iota \in H} \mathbf{E}_{\iota}\) ，其中 H 是 I 的一个有限
\end{enumerate}

子集，且 B 在每个 \(E_{\iota}\) （\(\iota \in H\) ）上的投影是有界的。由此推出（对 K = R 或 C）：若每个 \(E_{\iota}\) 都是拟完备空间，则赋以 \(T_{0}\) 的 E 是拟完备的。

\begin{enumerate}
\def\labelenumi{\arabic{enumi})}
\setcounter{enumi}{10}
\tightlist
\item
  设 \(\mathbf{E}\) 是带非离散赋值的域 \(\mathbf{K}\) 上的一个拓扑向量空间。
\end{enumerate}

\begin{enumerate}
\def\labelenumi{\alph{enumi})}
\item
  为使 E 的均衡子集 A 吸收 E 的每个有界子集（III, 第 37 页, 习题 1），只需 A 吸收 E 中每个趋于 0 的序列 \((x_{n})\) 的点所成的集即可。此时称 A 是吸有界集的。
\item
  设 \(u\) 是从 \(E\) 到一个拓扑向量空间 \(F\) （\(K\) 上的）中的线性映射。\(E\) 的每个有界子集在 \(u\) 下的像在 \(F\) 中有界，当且仅当对每个点列 \((x_{n})\) （\(E\) 中的、趋于 0 的），序列 \((u(x_{n}))\) 在 \(F\) 中有界。
\item
  假设 \(\mathbf{E}\) 可度量化。证明：\(\mathbf{E}\) 的每个吸有界集的子集都是 \(\mathbf{E}\) 中 0 的一个邻域。由此推出：若 \(u\) 是从 \(\mathbf{E}\) 到一个拓扑向量空间 \(\mathbf{F}\) （\(\mathbf{K}\) 上的）中的线性映射，它把 \(\mathbf{E}\) 中每个收敛于 0 的序列变为 \(\mathbf{F}\) 中的有界序列，则 \(u\) 在 \(\mathbf{E}\) 上连续。
\end{enumerate}

\begin{enumerate}
\def\labelenumi{\arabic{enumi})}
\setcounter{enumi}{11}
\item
  设 E 是带非离散赋值的域 K 上的 Hausdorff 拓扑向量空间，F 是 K 上的可度量化向量空间。若 u 是从 E 到 F 中的连续线性映射，使得对 F 的每个有界子集 B，\(u^{-1}(\mathbf{B})\) 在 E 中有界，证明 u 是从 E 到 F 的一个子空间上的同构。
\item
  设 \(I\) 是无穷集，\((E_{\iota})_{\iota\in I}\) 是一族局部凸空间，其中没有一个是 0。设 \(f\) 是从 \(E = \prod_{\iota\in I}E_{\iota}\) 到 Banach 空间 \(F\) 中的一个线性映射。证明：若 E 的每个有界子集在 f 下的像都是 F 的有界子集，则存在 I 的有限子集 H，使得对每个 \(\iota \notin H\) ，f 在 \(E_{\iota}\) （视为 E 的子空间）上的限制是零。（用反证法：否则可以构造 E 中的有界序列 \((x_{n})\) ，使其在 f 下的像在 F 中无界。）
\item
  证明：若可度量化局部凸空间 E 的拓扑不能由单个范数定义，则 E 的典范有界结构不存在可数基（利用 III, 第 38 页, 习题 5，证明否则 E 中将存在一个有界的吸有界集的集合（III, 第 39 页, 习题 11），并利用 III, 第 39 页, 习题 11, c) 完成论证）。
\item
  在 R 上的 Hausdorff 拓扑向量空间 E 中，设 A 是紧凸集，B 是闭的凸有界集。证明：并 \(A \cup B\) 的凸包 C 是闭集。（考虑 C 的闭包中不属于 A 的一点 z，并化归到 z = 0 的情形。注意存在 0 的一个邻域 V 和一个数 \(\alpha < 1\) ，使得关系式 \(0 \leqslant \lambda \leqslant 1\) 、\(x \in A\) 、\(y \in B\) 、\(\lambda x + (1 - \lambda)y \in V\) 蕴含 \(\lambda \leqslant \alpha\) 。然后，对 0 的每个邻域 W，考虑三元组 \((\lambda, x, y)\) （满足 \(\lambda x + (1 - \lambda)y \in W\) 、\(0 \leqslant \lambda \leqslant 1\) 、\(x \in A\) 、\(y \in B\) ）全体所成的集。）
\item
  设 \(\mathbf{E}\) 是满足第一可数公理的可度量化局部凸空间，且其完备化 \(\hat{\mathbf{E}}\) 是满足第一可数公理的 Fréchet 空间。
\end{enumerate}

证明：\(\hat{E}\) 的每个有界子集 B 都含于 \(\hat{E}\) 的某个有界子集的闭包中。（化归到 B 可数、并排成序列 \((x_{n})\) 的情形；另一方面，设 \((p_{n})\) 是定义 \(\hat{E}\) 的拓扑的半范数的递增序列；对每个整数 n，考虑 E 的点列 \((y_{nk})_{k\geqslant1}\) ，它收敛于 \(x_{n}\) ，且满足 \(p_{n}(x_{n}-y_{nk})\leqslant1\) （对一切 \(k\geqslant1\)）。）

\begin{enumerate}
\def\labelenumi{\arabic{enumi})}
\setcounter{enumi}{16}
\tightlist
\item
  设 \(\mathbf{A}\) 是 \(\geqslant 1\) 的、从 \(\mathbf{N}\) 到 \(\mathbf{N}\) 中的递增映射全体所成的集；对每个 \(\alpha \in \mathbf{A}\) ，用 \(\mathbf{B}_{\alpha}\) 表示由所有点 \(z = (z_n) \in \mathbf{R}^{\mathbf{N}}\) （它们满足 \(|z_n| \leqslant \alpha(n)\) ，对一切 \(n \in \mathbf{N}\) ）组成的集。
\end{enumerate}

\begin{enumerate}
\def\labelenumi{\alph{enumi})}
\item
  证明：诸集 \(\mathbf{B}_{\alpha}\) 构成空间 \(\mathbf{R}^{\mathbf{N}}\) 的全体有界子集所成的有界结构的一个基。b) 对每个 \(\alpha \in \mathbf{A}\) ，集 \(\mathbf{RB}_{\alpha}\) 是 \(\mathbf{R}^{\mathbf{N}}\) 的向量子空间，它异于 \(\mathbf{R}^{\mathbf{N}}\) 且在 \(\mathbf{R}^{\mathbf{N}}\) 中稠密；因此存在一个线性型 \(f_{\alpha} \neq 0\) （不连续的，在 \(\mathbf{R}^{\mathbf{N}}\) 上），使得 \(f_{\alpha}(z) = 0\) 对一切 \(z \in \mathbf{B}_{\alpha}\) 成立。
\item
  设 \(\mathbf{E}\) 是由所有映射 \(g: \alpha \mapsto (g_n(\alpha)) \in \mathbf{R}^{\mathbf{N}}\) （从 \(\mathbf{A}\) 到 \(\mathbf{R}^{\mathbf{N}}\) 中，且对一切 \(n \in \mathbf{N}\) 使和 \(p_n(g) = \sum_{\alpha} |g_n(\alpha)|\) 为有限）全体所成的向量空间。证明诸 \(p_n\) 是半范数
\end{enumerate}

它们在 E 上定义出一个 Fréchet 空间的拓扑。

\(d\) ) 设 \(\mathbf{H}\) 是由所有 \(h\in \mathbf{E}\) （满足 \(h(\alpha)\in \mathbf{RB}_2\) ，对一切 \(\alpha \in \mathbf{A}\)）组成的集；证明 \(\mathbf{H}\) 是 \(\mathbf{E}\) 的一个处处稠密的向量子空间（注意：若 \(h\in \mathbf{E}\) 满足 \(h(\alpha) = 0\) （除 \(\alpha \in \mathbf{A}\) 的有限多个值外），则它属于 \(\mathbf{H}\) ）。

\begin{enumerate}
\def\labelenumi{\alph{enumi})}
\setcounter{enumi}{4}
\tightlist
\item
  设 \(\mathrm{E}_0 \subset \mathrm{E}\) 是 \(\mathrm{E}\) 中由所有 \(g \in \mathrm{E}\) （满足 \(\sum_{\alpha} |f_{\alpha}(g(\alpha))| < +\infty\) ）组成的向量子空间；
\end{enumerate}

此时映射 \(u: g \mapsto (f_{\alpha}(g(\alpha)))_{\alpha \in \mathbf{A}}\) 是从 \(\mathrm{E}_0\) 到 Banach 空间 \(\mathrm{F} = \ell^1(\mathrm{A})\) 中的一个线性映射（I, 第 4 页）。证明 \(u(\mathrm{E}_0)\) 在 \(\mathrm{F}\) 中处处稠密（注意对每个有限子集 \(\mathrm{J}\) （\(\mathrm{A}\) 的），存在 \(g \in \mathrm{E}_0\) ，使得 \(g(\alpha) = 0\) （对一切 \(\alpha \in \mathrm{A} - \mathrm{J}\)），且诸 \(f_{\alpha}(g(\alpha))\) （对 \(\alpha \in \mathrm{J}\)）取 \(\mathbf{R}\) 中任意值）。证明 \(u^{-1}(0)\) 在 \(\mathrm{E}_0\) 中处处稠密（利用 \(d\) ）。最后证明：对每个有界子集 \(C\) （\(E\) 的），存在 \(\alpha_0 \in A\) ，使得 \(f_{\alpha_0}(g(\alpha_0)) = 0\) 对一切 \(g \in C \cap E_0\) 成立，并由此推出 \(u(C \cap E_0)\) 在 \(F\) 中的闭包不是 \(F\) 中 0 的邻域。

\begin{enumerate}
\def\labelenumi{\alph{enumi})}
\setcounter{enumi}{5}
\tightlist
\item
  设 G 是 u 在 \(E_{0} \times F\) 中的图像，它是 Fréchet 空间 \(E \times F\) 的一个向量子空间。证明 G 在 \(E \times F\) 中处处稠密（注意对每个 \(x \in E_{0}\) ，\(x + u^{-1}(0)\) 在 E 中稠密）。然而证明：对 G 的每个有界子集 M，闭包 \(\overline{M}\) （M 在 \(E \times F\) 中的）不包含有界集 \(\{0\} \times U\) （\(E \times F\) 中的），其中 U 是 F 中的单位球（若 \(N = \text{pr}_{1}(M)\) ，注意由 e)，\(\overline{u(N)}\) 不可能包含 U）。
\end{enumerate}

\begin{enumerate}
\def\labelenumi{\arabic{enumi})}
\setcounter{enumi}{17}
\tightlist
\item
  在 Banach 空间 \(\ell^1 (\mathbf{N})\) （I, 第 4 页）中，设 \(e_m\) 为这样的序列 \((\delta_{mn})_{n\geq 0}\) ：\(\delta_{mn} = 0\) （对 \(m\neq n\)）且 \(\delta_{nn} = 1\) 。定义一个从 \(\ell^1 (\mathbf{N})\) 到 \(\mathbf{R}\) 中的连续映射，它把由诸 \(e_n\) 组成的有界序列变为 \(\mathbf{R}\) 的一个无界子集（利用 Urysohn 定理（GT, IX, §4, No.~2, 定理 2））。
\end{enumerate}

